% Chapter styles; all are suitable for black-and-white printing.
% Copyright 2026 Adib Hasan. LPPL 1.3c or later.
\documentclass[oneside,openany]{bookbangla}
\setbanglamainfont{Kalpurush}
\title{অধ্যায়ের নকশা}
\newcommand\samplechapter{%
  \chapter{গণনার শুরু}
  \section{পায়রার খোপ নীতি}
  একটি পরিচিত ধারণা দিয়ে শুরু করি। বস্তু ও খোপ ঠিকভাবে চিহ্নিত করতে
  পারলেই অনেক গণনার সমস্যার সমাধান সহজ হয়ে যায়।
  \begin{example}
    আটজন বন্ধুর মধ্যে অন্তত দুজনের জন্ম একই বারে কেন?
  \end{example}
  \begin{solution}সপ্তাহে সাতটি বার আছে।\end{solution}}
\begin{document}
\mainmatter
\chapterstyle{classic}
\samplechapter
\begin{center}\verb|\chapterstyle{classic}|\end{center}
\chapterstyle{default}
\samplechapter
\begin{center}\verb|\chapterstyle{default}|\end{center}
\chapterstyle{modern}
\samplechapter
\begin{center}\verb|\chapterstyle{modern}|\end{center}
\chapterstyle{minimal}
\samplechapter
\begin{center}\verb|\chapterstyle{minimal}|\end{center}
\chapterstyle{ruled}
\samplechapter
\begin{center}\verb|\chapterstyle{ruled}|\end{center}
\end{document}
